\section{Problem Formulation}

We formulate emotional preservation and prosodic realignment during low-resource language adaptation as a fully observable Markov Decision Process (MDP). Operating within an autoregressive Text-to-Semantic (T2S) environment, the RL agent (a language model policy $\pi_\theta$) observes a state $s$ encompassing available conditioning information: an input text sequence, a speaker timbre embedding, a duration constraint embedding, and a target global emotion embedding. The action $a$ is the generation of a discrete semantic token sequence, which is subsequently converted into an acoustic waveform. Upon generating a complete sequence, the RL learner receives a multi-objective scalar reward signal $R(s, a)$ computed from the decoded speech. This composite feedback vector, formalized as $R = \lambda_1 R_{\text{emo-cls}} + \lambda_2 R_{\text{intensity}} + \lambda_3 R_{\text{emphasis}} + \lambda_4 R_{\text{ASR}}$, balances four task-specific scalar components: categorical emotion classification accuracy, global emotion intensity (distance to a neutral centroid in Valence-Arousal-Dominance space), local emphasis control (relative pitch and energy patterns), and an ASR-based content reward for semantic intelligibility.

The optimization problem seeks to update the policy parameters $\theta$ by maximizing the expected multi-objective reward using Group Relative Policy Optimization (GRPO). The maximization objective, $\mathbb{E}_{s \sim \mathcal{D}, a \sim \pi_\theta} [ R(s, a) \nabla_\theta \log \pi_\theta(a|s) ]$, is subject to a Kullback-Leibler (KL) divergence penalty constraint, $\beta \text{KL}(\pi_\theta(\cdot|s) \| p_{\text{adapt}}(\cdot|s))$. This constraint anchors the policy to a base model $p_{\text{adapt}}$ to prevent catastrophic deviation and preserve the pretrained high-resource expressive prior. Our formulation rests on two key assumptions: (1) an initial parameter-efficient adaptation stage successfully aligns the model to the target low-resource language while keeping most expressive knowledge intact, and (2) the discrete semantic token space produced by the T2S module is sufficiently expressive to retain and recover subtle emotional intensities and local prosody variations during final acoustic reconstruction.
